\begin{pdSolution}{fh:ex:fh-xfer-attenuation}
The proof sketch's move --- ``each attenuation is a subset operation in the Anchor algebra; subset composition is a subset'' --- is exactly right when $\mathrm{att}_A$ and $\mathrm{att}_B$ both narrow authority \emph{within one shared capability/resource universe}. It stops being a statement about the same resource space at the second attenuation: $\mathrm{att}_B$ does not narrow a set of operations on $A$'s test database, because $D_B$ has no test database to narrow into. It grants a \emph{new} authority, over a resource $D_A$ never named, that $D_B$ alone controls. Calling $C_B'$ a ``subset of $C_A$'' silently assumes a mapping between the two harbors' resource namespaces that nothing in the four-message ceremony (\S\ref{fh:sec:fh-xfer-ceremony}) actually constructs.

What minting $C_B'$ really is, in this case, is \textbf{token exchange and new authorization at $B$}, not attenuation of $C_A$. A correct $D_B$ must additionally establish: (a) \emph{$A$-side input eligibility} --- a signed semantic/resource mapping saying that $C_A$'s authority over $A$'s test database is the kind of thing that may be exchanged at all, recorded with its dependency on $A$'s issuing epoch root $R_A^{(e)}$ so a later revocation of $C_A$ is traceable to $C_B'$; and (b) \emph{$B$-side local authorization} --- $D_B$'s own policy minting a genuinely new child grant over its staging database, under $D_B$'s sovereignty (Design Invariant~\ref{fh:def:fh-sovereignty}, item 2), rather than treating $\mathrm{att}_B$ as if it were narrowing $A$'s grant.

One more gap rides along with this one and is worth naming here rather than assuming solved: Bob's harbor observing a failed acceptance check, or an inclusion proof against $R_B^{(e)}$, does not by itself establish that Alice's agent \emph{caused} the failure (\S\ref{fh:sec:fh-settle-proto}'s damage-claim step). Settlement evidence must bind the exact work-unit, the base and result trees, the mediated effects, the acceptance policy in force, and the verifier that checked it; semantic causation --- ``$\alpha$'s write is why the check failed'' --- remains a claim for the adjudicator named in \S\ref{fh:sec:fh-escrow}'s 2-of-3 instantiation to weigh, not something the inclusion proof settles on its own. Theorem~\ref{fh:lem:fh-att} is not false; it is scoped to the case its proof sketch actually covers, and this exercise is that scope made concrete.
\end{pdSolution}
\begin{pdSolution}{fh:ex:fh-revocation-finality}
All three configurations share one module; only the \texttt{Rollback} and \texttt{Epoched} constants change. The committed results (\path{proofs/relay/CardRevocation.md}, ``Results (TLC v1.8.0, Homebrew OpenJDK)''), quoted exactly:
\begin{lstlisting}
baseline:  Model checking completed. No error has been found.   (13 distinct states)
rollback:  Invariant RevocationFinal is violated.               (backup-restore re-accepts a revoked card)
epoch:     Model checking completed. No error has been found.   (76 distinct states)
\end{lstlisting}
\textbf{Baseline} (\texttt{Rollback=FALSE, Epoched=FALSE}, \path{CardRevocation_baseline.cfg}): no adversary action is available, so a card in \texttt{revoked} stays revoked and \texttt{RevocationFinal} holds trivially over 13 reachable states. \textbf{Rollback} (\texttt{Rollback=TRUE, Epoched=FALSE}, \path{CardRevocation_rollback.cfg}): \texttt{RollbackRestore} can drop a card from the DB's \texttt{revoked} set while \texttt{everRevoked} (ground truth) stays set; \texttt{Accept} then only checks \texttt{revoked}, so the restored card is accepted again and \texttt{badAccept} fires --- the invariant is violated \emph{by design}, as the negative control for the attack. \textbf{Epoch} (\texttt{Rollback=TRUE, Epoched=TRUE}, \path{CardRevocation_epoch.cfg}): the mitigation adds an external epoch that \texttt{Revoke} bumps and that a restore does \emph{not} roll back; \texttt{Accept} now additionally requires \texttt{cardEpoch[c] = epoch}, so a restored card's stale epoch disables acceptance even though \texttt{revoked} was rolled back, and the invariant holds again over the larger 76-state space (the mitigation adds reachable states without reintroducing the breach).

The rollback configuration tests whether a restored database can resurrect a revoked capability. A relay revokes card $c$ (an agent's leaked capability); the relay's database is later restored from a backup taken before the revocation; the restored \texttt{revoked} set no longer contains $c$; the relay's own \texttt{Accept} check, consulting only that set, admits a presentation of $c$ it had already promised to refuse. TLC's counterexample trace is exactly \texttt{Issue(c)} $\to$ \texttt{Revoke(c)} $\to$ \texttt{RollbackRestore(c)} $\to$ \texttt{Accept(c)}, four transitions, \texttt{badAccept} true at the end. The failure is the intended result because a baseline or an untested mitigation would give false comfort: it is precisely the negative control that lets the epoch configuration's clean pass mean something, the same discipline \S\ref{fh:sec:fh-sheaf}'s negative-control ProVerif line and mutation suite apply. This chapter's own per-epoch roots $R_X^{(e)}$ (\S\ref{fh:sec:fh-witness}) are the same design move at the federation layer: an external, monotonic anchor that a local rollback cannot touch is what makes ``only-issuer-revokes'' something a cross-harbor verifier can actually rely on, rather than merely something the issuer asserts.
\end{pdSolution}
\begin{pdSolution}{fh:ex:fh-consistency-radius-run}
By hand: $C_6$ has $R_{\mathrm{eff}}=5/6$, so $r=3\sqrt{1-5/6}=3/\sqrt6=1.224745$ [verified]. $P_6$'s relayed edge is a bridge (a path has no route between an edge's endpoints other than the edge itself), so $R_{\mathrm{eff}}=1$ and $r=3\sqrt{1-1}=0$ exactly [verified].

The script's own printed lines (\path{skills/harbor-results/scripts/sheaf_consistency_radius.py}, \texttt{section0}), quoted in the worked example above [internal, \texttt{sheaf\_consistency\_radius.py}]: the $C_6$ line's two computed quantities (the least-squares solve's $r$ and the closed form's \texttt{pred}) are asserted equal to within $10^{-9}$ of each other and within $10^{-6}$ of $1.224745$ (\texttt{check(abs(r - pred) < 1e-9 and abs(r - 1.224745) < 1e-6, \ldots)}) --- the numerical solve and the hand formula are checked against each other, not merely printed side by side. The $P_6$ line's assertion is different in kind: \texttt{check(rP < TOL\_SILENT and abs(ReffP - 1.0) < 1e-9, \ldots)} with \texttt{TOL\_SILENT}$=10^{-9}$, i.e.\ it asserts the residual is \emph{below a silence threshold}, not equal to a specific float --- the honest way to assert ``zero'' about a least-squares solve that carries ordinary floating-point roundoff rather than landing on the exact bit pattern for $0.0$.
\end{pdSolution}
\begin{pdSolution}{fh:ex:fh-two-liar-boundary}
$r$ is defined as the \emph{best-fit} residual: the smallest disagreement left over after choosing the global section $x$ and the free severed blocks $g_{\mathrm{sev}}$ that explain as much of the observed data as any choice can. It is not a per-liar quantity; it is a property of the whole observed cochain $g_K$. If two equivocators on a common cycle choose offsets $o_1$ and $o_2$ such that $o_1+o_2$ happens to be a coboundary --- i.e.\ the \emph{sum} of their two lies is indistinguishable from a single consistent state change --- then the best-fit global section can absorb $o_1+o_2$ entirely, leaving a near-zero residual, exactly the way a lone liar's offset is absorbed when it lies off a visible cycle (the cut-edge case above). Nothing in the definition of $r$ requires the offsets to be individually small for their \emph{sum} to be small or zero; it requires only that the sum land in the image of $\delta$, the space the completion can always explain away. This is why the single-equivocator hypothesis is decisive rather than a simplifying nicety: the moment a second colluding reporter is admitted on the same cycle, the detector's soundness guarantee (Theorem~\ref{fh:thm:sheaf-equivocation}'s ``if and only if'') no longer applies to the \emph{pair}, only to each alone under the counterfactual that the other is honest.

The compendium's own sweep states the shape of this exactly: a solo equivocator on an 8-cycle is bounded below at $r\ge|s|/\sqrt8$, while a coordinating coalition of two on the same cycle drives the same statistic to $r\approx6\times10^{-15}$ --- numerically the silence of a true coboundary, not a small-but-detectable residual [internal, \texttt{sheaf\_consistency\_radius.py}]. What could still catch the coalition is exactly what $r$ is not: per-report attribution rather than a single scalar. Signature-level forensics can ask whether the \emph{same} reporter's two prefixes are individually self-consistent (an equivocator still signs two incompatible things even if the aggregate looks clean); cross-checking against a third, independent visible cycle through only one of the two colluders can isolate that one's contribution if such a cycle exists; and widening $K$ --- comparing more edges directly rather than relying on relayed reports --- shrinks the coalition's room to choose a cancelling pair in the first place. None of this is a proof that coalitions are caught in general; \S\ref{fh:sec:fh-lim-cartel} names the cross-harbor cartel question as open for exactly this reason, and this exercise is that same boundary stated at the level of one theorem rather than the whole federation.
\end{pdSolution}
\begin{pdSolution}{fh:ex:fh-proverif-federation}
The four principals are the daemon's local DB ($D$, read-only), the KMS ($K$, read-only key witness), the email provider ($E$, read access to the user's mailbox), and the passphrase ($P$, out-of-band, held only in the user's head --- deliberately not modeled as compromisable, since \S7 does not defend against the user revealing their own passphrase). The single query the file states is \texttt{query attacker(account\_root)}: can the attacker ever derive the reconstructed vault. The model has no destructor for \texttt{combine4}, so the four shares can only be combined by the honest user's own process, never inverted by an attacker who merely holds some subset of them.

Only one scenario is actually run --- the top-level \texttt{process} compromises $D$, $E$, and $K$ simultaneously (three of four), leaving $P$ uncompromised. The file's own comment states why this subsumes the weaker cases: ``if the vault is safe with daemon+email+KMS public, it is also safe with any subset of those compromised.'' An attacker who additionally lacks one of $D$, $E$, or $K$ is strictly weaker than the modeled attacker, so a query that holds against the strongest three-of-four attacker holds against every subset of it for free; running the weaker scenarios separately would add nothing this run does not already dominate.

Pinned verdict (\path{proofs/bonded/federated/federated.run.log}), quoted exactly:
\begin{lstlisting}
Starting query not attacker(account_root[])
RESULT not attacker(account_root[]) is true.
--------------------------------------------------------------
Verification summary:

Query not attacker(account_root[]) is true.
\end{lstlisting}
What this does \emph{not} cover, stated as plainly as \S\ref{fh:sec:fh-xfer-mech} states its own draft status: all four principals compromised at once (explicitly out of scope --- the four-eyes attacker), and the \S7.5 Shamir 3-of-4 recovery path, which is the honest user's own reconstruction option, not an attacker capability. Neither this model nor the cross-realm sketch above substitutes for the other: one is a completed secrecy result about who must collude for a vault to fall, the other is this chapter's own in-progress model of cross-realm card authenticity.
\end{pdSolution}
\begin{pdSolution}{fh:ex:fh-correlated-equilibrium}
One model treats the two harbor principals as players, with their daemons executing the chosen policies. Actions include accept, defer, and reject a transfer; payoffs include completed-work value, execution cost, delay, and enforceable losses. Each principal observes its local history and the signed remote evidence actually delivered to it. The model must specify a prior over hidden state and a delivery process; roots merely present in a log are not observations by both players.

For a complete-information stage-game simplification, let $\mu$ be the proposed joint distribution of actions. Coarse correlated equilibrium requires, for every player $i$ and fixed alternative action $a_i'$,
\[
  \mathbb{E}_{a\sim\mu}[u_i(a)]
  \ge
  \mathbb{E}_{a\sim\mu}[u_i(a_i',a_{-i})].
\]
Here the alternative is chosen before seeing the recommendation. Correlated equilibrium also excludes profitable deviations conditioned on the player's recommendation~\pdcite{mega281}. With private types or additional public signals, the relevant information and conditional deviations must be modeled explicitly; the displayed stage-game condition is not a proof for the richer protocol.

Fresh roots do not establish mutual receipt, common knowledge, or obedience. A proposed common-view event needs an explicit delivery and acknowledgement model, and its resulting information structure still needs an incentive check. If required evidence is missing or stale, defer the dependent cross-harbor effect; continue only independently authorized local work and preserve pending obligations under the settlement rules. This is a conservative authorization policy, not a proved equilibrium fallback. The federation result remains open.
\end{pdSolution}
\begin{pdSolution}{fh:ex:fh-trustless-settlement}
An HTLC's atomicity rests on a single objective, third-party-checkable bit: does a preimage of a fixed hash appear on-chain before a deadline. Both parties' payoffs are a deterministic function of that one bit, so the same contract executes identically for anyone who can read the chain. The Federated Harbor's damage predicate (\S\ref{fh:sec:fh-settle-proto}'s claim step) is not a single objective bit revealed by one party: whether $\alpha$'s action at $B$ actually caused the acceptance-criteria failure is a semantic, evidence-weighing judgment over a work-unit's base and result trees, exactly the causation gap Exercise~\ref{fh:ex:fh-xfer-attenuation}'s solution names. A deterministic trustless contract must settle identically whenever its on-ledger inputs are identical; if two possible worlds produce the same signed evidence trail but differ in whether $\alpha$ was actually at fault, no deterministic on-ledger rule can be correct in both worlds at once. That is the shape of the impossibility this exercise is asked to sketch, not merely a practical gap: it would need to assume exactly two worlds with identical observable transcripts and provably different ground truth, and show every candidate deterministic settlement rule is wrong in at least one of them. Progress, as \S\ref{fh:sec:fh-lim-trustless} and \S\ref{fh:sec:fh-escrow}'s 2-of-3 candidate both note, means adding something outside pure on-ledger determinism --- an oracle, a TEE, an optimistic dispute window, or a bonded human/arbiter adjudicator --- each of which reintroduces a trust assumption rather than removing one. This chapter gives neither the construction nor the impossibility proof; it gives the escrow's conditional bound (Design Invariant~\ref{fh:thm:fh-escrow-bound}) as the honest fallback while the question stays open.
\end{pdSolution}
\begin{pdSolution}{fh:ex:fh-cartel-formation}
For $p\cdot d\cdot L>G$ to be satisfiable against a majority-but-not-total cartel, the audit mechanism needs an independent source of $p$ and $d$ that the cartel does not itself control: cross-harbor adversarial audits sampled from outside the colluding set, witness diversification so that a majority of witnesses being compromised still leaves an honest witness able to observe and attest to the true state, and principal-wide collateral so that $L$ scales with what the colluding \emph{principals} (not just their harbor identities) have at stake, per \S\ref{fh:sec:fh-lim-cartel}'s list. Concavity of repeated-game payoffs in reputation credit raises the effective $G$ a cartel forfeits by being caught once, which is what makes $p\cdot d\cdot L>G$ achievable at plausible $p$ rather than requiring $p$ near 1. All of this requires that at least one witness and one audit path lie genuinely outside the cartel's control --- the mechanism's leverage comes entirely from an outside observer existing.

That is exactly why no structural mechanism closes the all-witnesses case: if the cartel controls every oracle and witness in scope, there is no outside observer left to supply $p$ or $d$ at all, and the inequality has no independent terms to be satisfied with. This is not a gap this chapter failed to close; it is the same limit every audit-based mechanism in the literature hits, stated honestly rather than argued around. \S\ref{fh:sec:fh-failure-modes}'s centralization-gravity discussion is the operational form of the same fact: the fewer independent witnesses a federation actually has, the closer it sits to the all-controlled boundary where no deterrence inequality can help.
\end{pdSolution}
\begin{pdSolution}{fh:ex:fh-cold-start-admission}
Four measurements, each cheap to compute off the witness-logged admission records \S\ref{fh:sec:fh-admission} already produces because enrollment is permanent and public: \textbf{acceptance rate} (admitted harbors / sponsorship requests, over a trailing window); \textbf{sponsor concentration} (the Herfindahl index, or simply the top-sponsor's share, of successful sponsorships); \textbf{time-to-entry} (elapsed time from first sponsorship request to admission, for accepted applicants); and \textbf{false-rejection rate} (requests declined that a later, independent sponsor accepted without incident, as a proxy for rejections that were not actually risk-justified). A thin-but-open market shows a low acceptance rate \emph{and} low sponsor concentration \emph{and} a time-to-entry that is long but not systematically longer for any identifiable class of applicant. An invitation cartel shows the opposite signature on the second and fourth: concentration rising over time even as total admissions grow, and a persistent gap in acceptance rate or time-to-entry between applicants introduced by an incumbent sponsor and applicants without one. None of these four numbers is a proof of openness by itself --- a legitimately small, honest market can also show high concentration simply because there are few sponsors to begin with --- which is exactly why \S\ref{fh:sec:fh-lim-cold-start} calls the structural pushback engineering discipline rather than theorem: the measurement tells an operator when to worry, it does not certify the absence of the failure mode.
\end{pdSolution}
\begin{pdSolution}{fh:ex:fh-equivocation-speed}
They are different problems because they sit on different sides of the same boundary. The deployment-specific tail bound is an \emph{empirical} gap: Theorem~\ref{fh:thm:fh-conv} already gives an expected $\Theta(\log m)$ result under a stated model (a complete overlay, reliable synchronous rounds), and what is missing is fitting that model's parameters --- or a weaker model's --- against a real peer graph's loss rates, fan-out, and partition frequency, exactly as \S\ref{fh:sec:fh-lim-equiv}'s solution direction calls for instrumenting real peer graphs and fitting tail quantiles before pricing a bond from propagation risk. The sheaf-Laplacian question is a \emph{mathematical} gap: \S\ref{fh:sec:fh-sheaf}'s closing paragraph is explicit that no relaxation-time bound has been derived or checked, because the Hansen--Ghrist rate governs continuous-time \emph{linear} diffusion over real stalks, while anti-entropy gossip over signed append-only logs under a Byzantine equivocator is discrete, asynchronous, adversarial, and monotone-join rather than linear-averaging --- closing it means either proving the two convergence stories are secretly the same rate, or proving they are provably different and explaining why the sheaf spectral gap is the wrong handle. Even a complete answer to one gap would not touch the other: a perfectly fitted empirical tail bound says nothing about whether $\lambda_{\min}^{+}(\Delta_{\mathcal F})$ predicts it, and a proved reconciliation of the two convergence models would still need real peer-graph data to become a number an operator can price a bond against. The measurement needed before either can be closed is the same first step: real anti-entropy round traces from an actual multi-harbor deployment, since \S\ref{fh:sec:fh-lim-equiv}'s empirical gap needs them directly and any attempt to validate a sheaf-diffusion reconciliation against gossip behavior needs the same trace data to compare against.
\end{pdSolution}
\begin{pdSolution}{fh:ex:fh-synthesis}
What does not require consensus is the \emph{event history}: harbors never need to agree on a single global order of issuances, revocations, or transfers. Each harbor's epoch roots only extend (per-harbor monotonicity), each root is cross-witnessed rather than centrally ordered, and equivocation is detectable rather than prevented by agreement (\S\ref{fh:sec:fh-not}, \S\ref{fh:sec:fh-witness}). What the Cross-Harbor Bucket Partition result (Design Invariant~\ref{fh:prop:fh-cross-cons}) presupposes is narrower and different in kind: it is a model-checked invariant of one abstract transition system with a fixed set of buckets per bond, checked at a stated finite bound ($|F|=4$, bond depth 16, \S\ref{fh:sec:fh-tla-ext}). Apalache is not requiring the \emph{harbors} to agree on anything at runtime; it is verifying that the \emph{rules} governing one bond's bucket moves cannot, on paper, create or destroy value, for every reachable state up to the bound. Consensus among machines and exhaustive checking of a specification are different axes entirely --- the first is a runtime coordination requirement the federation deliberately avoids, the second is an offline proof technique applied to a design the federation still has to implement correctly and run at unbounded scale. Nothing here claims the deployed custody ledger enforces the bound at runtime; that is exactly Design Invariant~\ref{fh:thm:fh-escrow-bound}'s conditional custody-ledger assumption, stated separately.

Calling the Federated Harbor ``a project institution, not merely two machines sharing a repository'' means the paper's actual claim is broader than any one protocol: an authoritative event/work-unit graph, per-principal policy, local effect mediators, causal messages, evidence, roles, and explicit governance over conflicts, of which the transfer ceremony, revocation mesh, and settlement escrow are three essential pieces, not the whole structure. A roadmap or a stated intention is versioned forecast, not truth, and a contradiction detector (the witness log's equivocation check) produces evidence for a \emph{named principal or role} to decide against, rather than deciding by itself. The paper's own posture --- consensus only where it is cheap and structurally necessary (the small authoritative metadata core: epoch roots, admission records), federation everywhere else (the actual work artifacts) --- is the concrete answer to the framing question of \S\ref{fh:sec:fh-intro}: two machines needed more than a shared repository, and what they needed is exactly the institution this chapter specifies one layer of.
\end{pdSolution}
